%=============================================================================
% LITERATURE REVIEW: QUANTUM DECOHERENCE, ERROR CORRECTION,
% GRAVITATIONAL DECOHERENCE, AND DFD QUANTUM ADVANTAGES
%
% Supporting Patent #3: psi-Field Quantum Control and Error-Suppression
% Gary Thomas Alcock
%=============================================================================
\documentclass[11pt,letterpaper]{article}

\usepackage{amsmath,amssymb}
\usepackage{graphicx}
\usepackage{booktabs}
\usepackage{hyperref}
\usepackage{xcolor}
\usepackage{geometry}
\usepackage{enumitem}
\usepackage{longtable}

\geometry{margin=1in}

\title{Literature Review:\\
Quantum Decoherence, Error Suppression, Gravitational Decoherence,\\
and DFD $\psi$-Field Quantum Advantages\\[6pt]
\large Supporting Research for Patent \#3}

\author{Gary Thomas Alcock\\
\textit{Independent Researcher, Los Angeles, CA, USA}}

\date{\today}

\begin{document}
\maketitle
\tableofcontents
\newpage

%=============================================================================
\section{Current Decoherence Rates Across Qubit Platforms}
%=============================================================================

\subsection{Superconducting Qubits}

\subsubsection{Transmon Qubits}

The transmon qubit, introduced by Koch et al.\ (2007), operates in the
regime $E_J/E_C \gg 1$ to suppress charge noise sensitivity. Key
coherence milestones:

\begin{itemize}
\item \textbf{2011:} Paik et al.\ (Yale) demonstrated $T_1 \approx
70\,\mu$s in 3D transmon architecture, a $10\times$ improvement over
planar designs. The improvement was attributed to reduced dielectric
participation ratio in the larger 3D cavity mode volume.

\item \textbf{2021:} Place et al.\ (Princeton) reported $T_1 > 300\,\mu$s
using tantalum capacitor electrodes. The Ta films, grown epitaxially on
sapphire, showed dramatically reduced two-level system (TLS) defect
density compared to niobium. Notably, the coherence improvement
($\sim 4\times$) exceeded the measured TLS reduction ($\sim 2\times$),
suggesting additional mechanisms.

\item \textbf{2023:} Somoroff et al.\ demonstrated millisecond-scale
coherence in fluxonium qubits ($T_1 \approx 1.5$\,ms at the sweet spot).
Fluxonium achieves longer $T_1$ than transmon due to its smaller
transition matrix element to the environment.

\item \textbf{2024--2025:} Several groups report $T_1 > 500\,\mu$s for
transmons with improved surface treatments (HF etching, vacuum annealing).
IBM, Google, and Rigetti all converge on $T_1 \sim 200$--$500\,\mu$s as
the state of the art for production-grade transmons.
\end{itemize}

\paragraph{Dominant noise sources:}
\begin{itemize}
\item $1/f$ flux noise: amplitude $A_\Phi \approx 1$--$5\,\mu\Phi_0/\sqrt{\rm Hz}$
\item TLS defects: Lorentzian spectrum, frequency-dependent density
\item Quasiparticle tunneling: exponentially suppressed below $\sim 150$\,mK
\item Thermal photon noise: $\bar{n}_{\rm th} \sim 0.01$--$0.1$ at
20\,mK in typical readout cavities
\end{itemize}

\subsubsection{Fluxonium Qubits}

Fluxonium operates at the half-flux-quantum sweet spot where first-order
flux sensitivity vanishes. Key results:
\begin{itemize}
\item $T_1$ up to 1.5\,ms (Somoroff et al., 2023)
\item $T_2^{\rm echo}$ up to 1\,ms
\item Anomalous $T_2/T_1$ ratios up to $\sim 5$ observed in some devices
--- exceeding the theoretical limit $T_2 \leq 2T_1$ for pure dephasing
\end{itemize}

The anomalous $T_2/T_1 > 2$ is significant for DFD: it suggests an
additional coherence protection mechanism beyond standard noise models.
In the DFD framework, the sweet-spot condition (vanishing first-order flux
sensitivity) may coincide with a local extremum of the $\psi$-potential
coupling, providing an additional stabilization channel.

\subsubsection{Coherence Limits and Projections}

The consensus view is that superconducting qubit coherence is limited by:
\begin{enumerate}
\item Material interfaces (metal--substrate, metal--air, substrate--air)
\item TLS defects in Josephson junction barriers
\item Quasiparticle poisoning
\item Infrared radiation
\end{enumerate}

Projected limits with perfect material engineering: $T_1 \sim 10$\,ms
(limited by fundamental quasiparticle generation from cosmic rays and
radioactivity). This ceiling motivates the search for fundamentally new
coherence-protection mechanisms---exactly what $\psi$-field stabilization
provides.

\subsection{Trapped-Ion Qubits}

\subsubsection{Hyperfine Qubits}

\begin{itemize}
\item $^{171}$Yb$^+$: $T_2 > 10$\,min at the magnetic-field-insensitive
``clock'' transition (Harty et al., 2014). Residual decoherence from
second-order Zeeman shift and oscillator phase noise.
\item $^{43}$Ca$^+$: $T_2 > 50$\,s reported (Ballance et al., 2016).
\item $^9$Be$^+$: $T_2 \approx 1.5$\,s (Langer et al., NIST).
\end{itemize}

\subsubsection{Optical Qubits}

\begin{itemize}
\item $^{40}$Ca$^+$ $S_{1/2} \leftrightarrow D_{5/2}$: natural lifetime
$T_1 \approx 1.2$\,s. $T_2$ limited by laser linewidth to
$\sim 1$--$10$\,s with dynamical decoupling.
\item $^{88}$Sr$^+$: similar performance.
\end{itemize}

\paragraph{Dominant noise:} Magnetic field fluctuations ($1/f$ and
line-frequency harmonics), ion-trap voltage noise, laser phase noise.

\subsection{Photonic Qubits}

\subsubsection{Cavity-Based Encodings}

\begin{itemize}
\item Superconducting 3D cavities: $T_1 \approx 2$\,ms (Yale,
$Q \sim 10^{10}$)
\item Cat-state encodings: $T_{\rm logical} \approx 1$\,ms demonstrated
(Amazon/Yale, 2024)
\item Binomial codes: error-corrected photonic qubit with
$T_{\rm logical} > T_{\rm bare}$ demonstrated
\end{itemize}

\subsubsection{Optical Photons}
\begin{itemize}
\item Fiber loss: $\sim 0.2$\,dB/km ($T_1 \sim 10\,\mu$s for 1\,km
fiber at $c/n$)
\item Free-space: limited by mode matching and detection efficiency
\item Linear optical QC: no intrinsic decoherence, limited by gate
success probability and photon loss
\end{itemize}

\subsection{Solid-State Spin Qubits}

\begin{itemize}
\item NV centers in diamond: $T_2 \sim 1$\,ms (with DD); $T_2^* \sim
10\,\mu$s
\item Silicon quantum dots: $T_2^* \sim 20\,\mu$s; $T_2^{\rm echo}
\sim 28$\,ms (Veldhorst et al., 2014)
\item $^{28}$Si:P donors: $T_2 > 30$\,s in isotopically purified
silicon (Muhonen et al., 2014)
\end{itemize}

%=============================================================================
\section{Dynamical Decoupling: Theory and Limits}
%=============================================================================

\subsection{Fundamental Principle}

Dynamical decoupling applies sequences of $\pi$-pulses to refocus
accumulated phase errors. The basic idea traces to the Hahn echo (1950):
a single $\pi$-pulse at $t/2$ refocuses static inhomogeneous broadening,
extending $T_2^*$ to $T_2^{\rm echo} \leq 2T_1$.

\subsection{Advanced DD Sequences}

\begin{itemize}
\item \textbf{CPMG} (Carr-Purcell-Meiboom-Gill): $N$ equally spaced
$\pi$-pulses. $T_2 \propto N^{2/3}$ for $1/f$ noise.
\item \textbf{UDD} (Uhrig Dynamical Decoupling): Non-uniform pulse
spacing optimized for specific noise spectra. Pulses at
$t_j = t\sin^2(\pi j/(2N+2))$. Optimal for soft (high-frequency
cutoff) noise spectra.
\item \textbf{CDD} (Concatenated DD): Recursive nesting of DD sequences.
Achieves $T_2 \propto N^{2^k/3}$ for $k$-th order concatenation.
\item \textbf{KDD} (Knill DD): Symmetrized pulses for robustness against
pulse errors.
\item \textbf{Adaptive DD:} Real-time optimization of pulse timing based
on measured noise spectrum (Sung et al., 2019).
\end{itemize}

\subsection{Theoretical Limits of DD}

DD faces fundamental limits:

\begin{enumerate}
\item \textbf{Pulse speed:} Finite-width pulses limit the maximum
refocusing rate. For pulse duration $\tau_p$: maximum DD frequency
$f_{\rm DD}^{\rm max} \approx 1/(2\tau_p)$. Noise above this frequency
is not suppressed.

\item \textbf{Pulse errors:} Each pulse introduces a small error
$\epsilon_p$. After $N$ pulses, accumulated error $\sim N\epsilon_p$.
For $\epsilon_p \sim 10^{-4}$, maximum useful $N \sim 10^4$.

\item \textbf{$T_1$ limit:} DD cannot extend $T_2$ beyond $2T_1$.
Energy relaxation is a $T_1$ process that $\pi_z$-pulses do not address.

\item \textbf{High-frequency noise:} DD is ineffective against noise at
frequencies $f > f_{\rm DD}^{\rm max}$. White noise and shot noise above
the DD bandwidth persist.

\item \textbf{Non-Gaussian noise:} DD theory assumes Gaussian noise
statistics. Non-Gaussian noise (e.g., telegraph noise from TLS switching)
requires modified treatment and may not be fully suppressed.
\end{enumerate}

\paragraph{Where DFD $\psi$-modulation goes further:}
\begin{itemize}
\item Addresses low-frequency noise directly (continuous modulation, not
discrete refocusing)
\item No accumulated pulse errors
\item Can suppress noise below DD's low-frequency cutoff
$f_{\rm min} \sim 1/t_{\rm gate}$
\item Multi-qubit correlated suppression through shared field
\item Compatible with DD as an additional layer
\end{itemize}

%=============================================================================
\section{Quantum Error Correction: Overhead and Limits}
%=============================================================================

\subsection{Surface Code}

The surface code (Kitaev, 1997; Fowler et al., 2012) is the leading
candidate for fault-tolerant QEC:
\begin{itemize}
\item Threshold: $p_{\rm th} \approx 1\%$ (with perfect measurements)
\item Overhead: $d^2$ physical qubits per logical qubit, where code
distance $d \sim \log(1/\epsilon_L)/\log(p_{\rm th}/p)$
\item At $p = 10^{-3}$: $d \approx 25$ for $\epsilon_L = 10^{-12}$,
requiring $\sim 625$ physical qubits per logical qubit
\item At $p = 10^{-5}$: $d \approx 9$, requiring $\sim 81$ physical
qubits per logical qubit
\end{itemize}

\subsection{Google Quantum AI Results (2023)}

The landmark Nature paper demonstrated the first scaling of logical error
rates with code distance in a surface code. Key findings:
\begin{itemize}
\item Physical error rate: $\sim 0.3\%$ (after optimization)
\item Distance-3 code: $\epsilon_L \approx 3.0\%$
\item Distance-5 code: $\epsilon_L \approx 2.9\%$ (slight improvement,
demonstrating below-threshold operation)
\item Extrapolation: $\sim 1000$ physical qubits needed per logical qubit
at current error rates for useful computation
\end{itemize}

\subsection{IBM Results (2024)}

IBM demonstrated error-mitigated circuits with up to 100+ qubits,
achieving useful computation through error mitigation rather than full
QEC. This approach works for near-term applications but does not scale
to arbitrary computation depth.

\subsection{QEC Overhead Problem}

The fundamental challenge: achieving $10^{-12}$ logical error rates
(required for Shor's algorithm on useful key sizes) requires:
\begin{itemize}
\item $\sim 20$\,million physical qubits at current error rates
\item $\sim 200{,}000$ physical qubits at $10\times$ improved error rates
\item $\sim 20{,}000$ physical qubits at $100\times$ improved error rates
\end{itemize}

$\psi$-field stabilization targeting $100$--$1000\times$ error rate
reduction would reduce the required qubit count to
$\sim 10{,}000$--$20{,}000$---within reach of near-term hardware.

%=============================================================================
\section{Gravitational Decoherence Theories}
%=============================================================================

\subsection{Penrose Proposal (1996)}

Penrose argued that the incompatibility between quantum superposition and
general relativity leads to a fundamental decoherence mechanism. For a
mass $m$ in superposition with separation $d$:
\begin{equation}
\tau_{\rm Penrose} \sim \frac{\hbar}{\Delta E_G}
\sim \frac{\hbar c^2 d}{G m^2}.
\end{equation}

For a $10^{-14}$\,kg mass with $d = 1\,\mu$m:
$\tau_{\rm Penrose} \sim 10^{-7}$\,s.

\subsection{Di\'osi Model (1987, 1989)}

Di\'osi proposed a continuous spontaneous localization model with
gravitational origin. The master equation adds a non-unitary term:
\begin{equation}
\frac{d\hat{\rho}}{dt} = -\frac{i}{\hbar}[\hat{H}, \hat{\rho}]
- \frac{G}{2\hbar}\int d^3x\,d^3y\,
\frac{[\hat{M}(\mathbf{x}), [\hat{M}(\mathbf{y}), \hat{\rho}]]}
{|\mathbf{x} - \mathbf{y}|},
\end{equation}
where $\hat{M}(\mathbf{x})$ is the mass density operator.

\subsection{Experimental Tests}

\subsubsection{Molecular Interferometry}

Fein et al.\ (Vienna, 2019): matter-wave interferometry with molecules
up to 25,000\,amu showing no gravitational decoherence. This constrains
the Penrose--Di\'osi model but does not definitively rule it out (mass
too low).

\subsubsection{Optomechanical Tests}

Vinante et al.\ (2017, 2020): cantilever decoherence measurements
finding excess noise consistent with the CSL model but at rates above
the Di\'osi prediction. Status: inconclusive.

\subsubsection{MAQRO/Space-Based Proposals}

The MAQRO proposal (Kaltenbaek et al.) would test gravitational
decoherence with $10^{-14}$\,kg particles in space, where isolation from
environmental noise is orders of magnitude better than on Earth.

\subsubsection{Tabletop Gravity-Entanglement (BMV)}

The Bose-Marletto-Vedral (BMV) proposal: create gravitational
entanglement between two mesoscopic masses in superposition. Detection
would prove gravity is quantum. DFD prediction: the experiment will show
gravitational effects consistent with the classical $\psi$-field
(expectation-value sourcing), not quantum gravity.

\subsection{DFD Resolution}

DFD eliminates gravitational decoherence entirely:
\begin{itemize}
\item No operator-valued geometry $\Rightarrow$ no Penrose paradox
\item Single classical $\psi$ field sourced by $\langle\hat{\rho}\rangle$
\item Convex-sum property: $\psi = |a|^2\psi_L + |b|^2\psi_R$
\item Well-posedness theorem: unique solution exists
\item Prediction: all gravitational decoherence experiments will find
null results (beyond environmental noise)
\end{itemize}

%=============================================================================
\section{Anomalous Coherence Observations Potentially Explained by DFD}
%=============================================================================

\subsection{Excess Coherence in 3D Transmon Architecture}

When Paik et al.\ (2011) placed a transmon in a 3D aluminum cavity, $T_1$
improved by $\sim 10\times$ relative to the best planar designs. The
standard explanation (reduced surface participation ratio) quantitatively
accounts for approximately $5\times$. The remaining factor of $\sim 2$
has been attributed to ``reduced spurious modes'' but without detailed
accounting.

\textbf{DFD interpretation:} The 3D cavity's EM mode structure creates a
more spatially uniform energy density distribution at the qubit location.
In DFD, this corresponds to a flatter $\psi$-landscape with smaller
$|\nabla\psi|$, reducing the kinetic coupling term in $\hat{H}_\psi$
(Eq.~(6) of the main paper).

\subsection{Tantalum ``Excess Improvement''}

Place et al.\ (2021) observed a $4\times$ improvement in $T_1$ when
switching from Nb to Ta electrodes, while TLS measurements showed only a
$\sim 2\times$ reduction in TLS density. The factor-of-2 discrepancy
remains unexplained.

\textbf{DFD interpretation:} Tantalum's different crystal structure
(BCC vs.\ Nb's BCC but with different lattice constant and electron
density) produces a different local $\psi$-coupling at the
metal--dielectric interface. The body-centered cubic Ta lattice, with its
higher density ($16.7$\,g/cm$^3$ vs.\ Nb's $8.6$\,g/cm$^3$), creates
a different $\psi$-gradient profile that partially stabilizes the qubit.

\subsection{Fluxonium $T_2/T_1 > 2$ Anomaly}

Several fluxonium devices show $T_2$ exceeding $2T_1$, which violates the
standard relation $T_2 \leq 2T_1$ valid for systems with pure dephasing
plus energy relaxation.

\textbf{DFD interpretation:} At the half-flux sweet spot, the vanishing
first-order flux sensitivity creates a condition where the $\psi$-field's
potential coupling (proportional to $\delta\psi$) provides a restoring
force against dephasing. This is analogous to a decoherence-free subspace
but arising from the physical $\psi$-landscape rather than from symmetry
of the system-bath coupling.

\subsection{Cavity-Protection Beyond Purcell}

Multiple groups observe that placing qubits inside high-Q cavities
provides coherence protection exceeding Purcell-filter predictions. The
``cavity protection effect'' is attributed to modified vacuum fluctuations,
but quantitative agreement with theory is imperfect.

\textbf{DFD interpretation:} The cavity's standing-wave EM field creates
a structured energy density landscape. At the field antinode, the local
$\psi$-value is higher, creating a potential well that stabilizes the
qubit phase. This is the physical-substrate mechanism underlying the
empirical cavity protection effect.

\subsection{Photon-Number-Dependent Coherence}

In the strong dispersive regime of circuit QED, qubit coherence varies
non-monotonically with cavity photon number $\bar{n}$. At specific
$\bar{n}$ values, $T_2$ peaks.

\textbf{DFD interpretation:} Each photon number state corresponds to a
different EM energy density, hence a different local $\psi$. At specific
$\bar{n}$, the $\psi$-potential creates a local extremum that provides
enhanced phase stability.

\subsection{Silicon Spin Qubit Record}

Muhonen et al.\ (2014) achieved $T_2 > 30$\,s for a $^{28}$Si:P donor
spin qubit in isotopically purified silicon. This is $\sim 10^3\times$
longer than in natural silicon.

\textbf{DFD interpretation:} Isotopic purification removes $^{29}$Si
nuclear spins (the dominant noise source). But it also creates a more
uniform mass density distribution, reducing $\psi$-gradients in the
crystal. The DFD prediction is that the coherence improvement should
slightly exceed the magnetic noise reduction---testable by precise
comparison of measured vs.\ predicted improvement factors.

%=============================================================================
\section{Refractive Index Modulation Effects on Quantum Coherence}
%=============================================================================

\subsection{Theoretical Background}

The connection between refractive index and quantum coherence is
well-established in optics:
\begin{itemize}
\item Slow light in EIT media enhances quantum memory times
\item Photonic band gaps in PhC structures modify vacuum fluctuations
\item Metamaterials with $n < 1$ or $n < 0$ modify photon dispersion
\end{itemize}

In DFD, $n = e^\psi$ is the fundamental field. Modulating $n$ modulates
the local phase velocity $c_1 = c/n$, providing direct control over
quantum phase evolution.

\subsection{Slow-Light Quantum Memories}

EIT-based quantum memories achieve storage times of $\sim 1$\,minute
(Dudin et al., 2013; Heinze et al., 2013) by reducing group velocity to
$\sim 10$\,m/s ($n_{\rm group} \sim 10^7$). The stored photon's
decoherence is determined by the spin coherence of the atomic ensemble.

\textbf{DFD connection:} The group index $n_g = n + \omega\,dn/d\omega$
in DFD includes a $\psi$-dependent contribution. Slow-light storage
corresponds to operation in a region of large $d\psi/d\omega$, which in
the DFD framework is a manifestation of dispersive $\psi$-coupling.

\subsection{Photonic Crystal Cavities}

Photonic crystal cavities achieve $Q > 10^7$ by engineering the
dielectric (refractive index) landscape. The EM energy density in these
structures creates a corresponding $\psi$-landscape that, in the DFD
framework, contributes to the total coherence environment.

\subsection{Metamaterial Substrates}

Superconducting metamaterials (arrays of SQUIDs, resonators, or
transmission lines) can create effective refractive indices from $n < 0$
(left-handed media) to $n \gg 1$. These structures are the natural
platform for $\psi$-field engineering.

%=============================================================================
\section{EM Field Environments That Improve Quantum Coherence}
%=============================================================================

\subsection{Decoherence-Free Subspaces}

Zanardi and Rasetti (1997) showed that certain subspaces of the
system Hilbert space are immune to collective decoherence. DFS protection
requires the noise to couple identically to all qubits---a condition
naturally satisfied by the shared $\psi$-envelope.

\subsection{Quantum Zeno Effect via Continuous Measurement}

Frequent measurement (or continuous monitoring) can suppress decoherence
by the quantum Zeno effect. The $\psi$-feedback loop implements an
analog version: continuous monitoring and correction rather than
projective measurement.

\subsection{Reservoir Engineering}

Poyatos et al.\ (1996) and subsequent work showed that engineered
dissipation can stabilize quantum states. The $\psi$-field provides a
novel form of reservoir engineering: the ``reservoir'' is the classical
scalar field, which can be externally controlled to drive the system
toward the desired quantum state.

\subsection{Floquet Engineering}

Periodic driving of quantum systems (Floquet engineering) can create
effective Hamiltonians with desired properties. The periodic $\psi$-modulation
is a form of Floquet engineering at the field level, creating an effective
Hamiltonian with reduced noise sensitivity.

%=============================================================================
\section{Quantum Error Correction Overhead Limits}
%=============================================================================

\subsection{Theoretical Minimum Overhead}

The quantum Hamming bound gives the minimum number of physical qubits
per logical qubit for a code correcting $t$ errors:
\begin{equation}
n \geq \frac{(3t)!}{t!(2t)!}
\end{equation}
For $t = 1$: $n \geq 5$. The 5-qubit code achieves this bound.

\subsection{Practical Overhead with Topological Codes}

Surface code overhead: $n \sim d^2$ where $d \sim O(\log(1/\epsilon_L))$.
Color code: similar scaling but with transversal gates.
Concatenated codes: $n \sim (\log(1/\epsilon_L))^{\log_2 r}$ where $r$ is
the code block size.

\subsection{Impact of Physical Error Rate}

The code distance $d$ needed for target $\epsilon_L$ scales as:
\begin{equation}
d \sim \frac{\ln(C/\epsilon_L)}{2\ln(p_{\rm th}/p)}
\end{equation}
where $C \sim 0.1$. Reducing $p$ by $100\times$ (from $10^{-3}$ to
$10^{-5}$) reduces $d$ by $\sim 3\times$ and $n$ by $\sim 9\times$.

This is the quantitative basis for the $\psi$-stabilization advantage:
$100\times$ coherence improvement translates to $\sim 9\times$ QEC
overhead reduction.

%=============================================================================
\section{Summary of DFD Quantum Advantages}
%=============================================================================

\begin{longtable}{p{4.5cm}p{4cm}p{4cm}}
\toprule
\textbf{Challenge} & \textbf{Standard Approach} & \textbf{DFD $\psi$-Field
Advantage} \\
\midrule
\endhead
Gravitational decoherence & Penrose collapse limits coherence &
Eliminated: single $\psi$ field, no branching \\[6pt]
$1/f$ flux noise & DD ($T_2 \propto N^{2/3}$) & Continuous cancellation
($T_2 \propto 1/(1-\eta_\psi)$) \\[6pt]
TLS defect noise & Material engineering & $\psi$-notch filter at defect
frequencies \\[6pt]
Multi-qubit correlated noise & Individual DD per qubit & Shared
$\psi$-envelope suppresses correlations \\[6pt]
QEC overhead ($\sim 10^3$ physical/logical) & Better materials,
gates & $100\times$ error rate reduction $\to$ $9\times$ overhead
reduction \\[6pt]
DD pulse errors & Composite pulses (KDD) & No pulses; continuous
field \\[6pt]
$T_1$ limit & Purcell filters, materials & $\psi$ addresses dephasing
channel independently \\[6pt]
Non-Gaussian noise (TLS) & Adaptive DD & Targeted $\psi$-cancellation
at TLS frequencies \\[6pt]
Anomalous coherence effects & Unexplained & Natural explanation via
$\psi$-coupling \\[6pt]
Scalability of error suppression & $O(N_q \cdot N_{\rm pulse})$ &
Single shared $\psi$-field \\
\bottomrule
\end{longtable}

%=============================================================================
\section{Key References and Experimental Targets}
%=============================================================================

\subsection{Critical Experiments for DFD Validation}

\begin{enumerate}
\item \textbf{Cavity--atom LPI test} (Alcock, 2025): $\xi \simeq 1$ in
DFD vs.\ $\xi = 0$ in GR. This is the gateway experiment. Signal:
$\Delta R/R \approx 1.1 \times 10^{-14}$ per 100\,m.

\item \textbf{Position-dependent coherence in 3D cavities:} Measure
transmon $T_2$ at different positions within a 3D cavity. DFD predicts
variation correlated with EM energy density profile.

\item \textbf{$\bar{n}$-dependent coherence mapping:} Systematically
measure $T_2$ vs.\ cavity photon number in the strong dispersive regime.
DFD predicts non-monotonic behavior with peaks at specific $\bar{n}$.

\item \textbf{Gravitational decoherence null test:} Any experiment
testing Penrose--Di\'osi collapse (MAQRO, tabletop optomechanics). DFD
predicts null result.

\item \textbf{$\psi$-substrate prototype:} Fabricate a resonator array
and demonstrate qubit frequency shift correlated with resonator photon
number --- first direct measurement of $g_\psi$.
\end{enumerate}

\subsection{Key Literature}

The following works are most relevant to Patent \#3:

\paragraph{Decoherence and noise:}
Koch et al.\ (2007): transmon design.
Bylander et al.\ (2011): noise spectroscopy via DD.
Slichter et al.\ (2012): measurement-induced transitions.
Klimov et al.\ (2018): fluctuations of energy-relaxation times.
Burnett et al.\ (2019): TLS decoherence.

\paragraph{Coherence records:}
Place et al.\ (2021): tantalum transmon.
Somoroff et al.\ (2023): fluxonium millisecond coherence.
Wang et al.\ (2022): long-lived transmon.
Muhonen et al.\ (2014): silicon donor spin.
Harty et al.\ (2014): trapped-ion record.

\paragraph{Error correction:}
Fowler et al.\ (2012): surface code.
Google Quantum AI (2023): below-threshold surface code.
Ofek et al.\ (2016): QEC extending cavity lifetime.
Sivak et al.\ (2023): real-time QEC.

\paragraph{Gravitational decoherence:}
Penrose (1996): gravity-induced collapse.
Di\'osi (1987, 1989): continuous localization.
Fein et al.\ (2019): molecular interferometry.
Donadi et al.\ (2021): underground collapse test.
Bose et al.\ (2017): gravitational entanglement.
Marletto and Vedral (2017): gravitational entanglement.
Alcock (2025): DFD Penrose resolution.

\paragraph{Dynamical decoupling:}
Viola et al.\ (1999): DD theory.
Uhrig (2007): optimal DD.
Biercuk et al.\ (2009): DD in trapped ions.
de Lange et al.\ (2010): DD in NV centers.

%=============================================================================
\section{Conclusions}
%=============================================================================

This literature review establishes that:

\begin{enumerate}
\item Quantum coherence times have improved dramatically over two decades
but remain orders of magnitude short of fault-tolerant requirements
without massive QEC overhead.

\item Dynamical decoupling has reached theoretical limits imposed by
pulse errors, finite pulse width, and the $T_1$ ceiling.

\item Gravitational decoherence (Penrose--Di\'osi) represents a
potentially fundamental limit in the GR framework that DFD entirely
eliminates.

\item Several anomalous coherence observations (excess improvements in
3D transmons, Ta ``bonus,'' fluxonium $T_2/T_1 > 2$, cavity protection
beyond Purcell) are naturally explained by DFD $\psi$-coupling effects.

\item The QEC overhead problem is a direct function of physical error
rate: $100\times$ improvement in coherence translates to $\sim 9\times$
reduction in physical qubit overhead.

\item DFD $\psi$-field stabilization provides a qualitatively new
mechanism---continuous, analog, correlated, multi-qubit noise suppression
at the physical substrate level---that complements all existing approaches.
\end{enumerate}

The experimental path forward is clear: verify the DFD $\psi$-coupling
through the LPI test, then demonstrate coherence enhancement in a
$\psi$-substrate prototype.

\end{document}
